% =============================================================================
% Draft mục 3.3.5 — Liquidity Service
% =============================================================================

\subsection{Liquidity Service}
\label{sec:liquidity-service}

Một matching engine chỉ thật sự hữu dụng khi luôn có đủ thanh khoản
để tạo thành cycle. Trong thực tế triển khai trên testnet (và trong
giai đoạn đầu của một bridge mainnet), số lượng intent đến từ người
dùng cuối thường thưa, dẫn đến việc nhiều intent phải chờ rất lâu
trước khi tìm được đối tác. \emph{Liquidity Service} của OceanLink
đóng vai trò bù đắp: tự động bơm các \emph{lệnh thanh khoản (LP order)}
do nhà cung cấp thanh khoản phát hành, dưới một cấu trúc kích thước
được thiết kế sao cho mọi nhu cầu của người dùng đều có thể được
khớp.

\subsubsection*{Bài toán: phủ thanh khoản tất định}

Đặt $A_{\max}$ là giới hạn trên của số lượng USDC mà một intent
người dùng có thể yêu cầu trong scope của hệ thống (mặc định
$A_{\max} = 10\,000$ USDC). Mục tiêu của Liquidity Service là đảm
bảo: \emph{với mọi giá trị $a \in [1, A_{\max}]$, tồn tại một tổ hợp
LP order có thể khớp tổng số $a$ trên đúng cặp chuỗi yêu cầu, dùng
ít LP order nhất có thể}.

\subsubsection*{Cấu trúc luỹ thừa của 2}

Liquidity Service gieo các LP order với mệnh giá là luỹ thừa của 2,
tương tự cách tiền giấy được phát hành theo các mệnh giá rời rạc:

\[
  \mathcal{D} = \{ 2^0, 2^1, 2^2, \dots, 2^{13} \}
              = \{ 1, 2, 4, 8, \dots, 8192 \}.
\]

Tổng các phần tử là $2^{14} - 1 = 16\,383 > A_{\max}$. Theo biểu diễn
nhị phân của một số tự nhiên, mọi $a \in [1, 16\,383]$ đều bằng tổng
một tập con của $\mathcal{D}$, và đây là tập con \emph{duy nhất}. Hệ
quả: matching algorithm chỉ cần chọn (greedy hoặc DP) các LP order
có mệnh giá $\le a$ để phủ chính xác lượng $a$ — số lượng order phải
dùng nhiều nhất là $\log_2 A_{\max} \approx 14$. Đây là phương pháp
chuẩn trong các hệ thống coin selection (ví dụ Bitcoin) và ATM cash
dispensing.

So với cách phát hành LP order với mệnh giá tuỳ ý hoặc đồng đều, cấu
trúc luỹ thừa của 2 vừa giảm số order cần duy trì
(\,$O(\log A_{\max})$ thay vì $O(A_{\max})$), vừa cho phép phủ
\emph{tất định}: không có giá trị nào trong khoảng $[1, A_{\max}]$
nằm ngoài khả năng phủ.

\subsubsection*{Phân bổ wallet LP và cặp chuỗi}

Hệ thống gắn ba ví LP với ba chuỗi nguồn:

\begin{itemize}
  \item Ví $B$ trên Ethereum Sepolia $\rightarrow$ tạo các LP order có đích
    là Base Sepolia và Arbitrum Sepolia.
  \item Ví $C$ trên Base Sepolia $\rightarrow$ tạo các LP order có đích là
    Ethereum Sepolia và Arbitrum Sepolia.
  \item Ví $D$ trên Arbitrum Sepolia $\rightarrow$ tạo các LP order có đích là
    Ethereum Sepolia và Base Sepolia.
\end{itemize}

Như vậy mỗi ví LP có hai cặp chuỗi nguồn–đích, và mỗi cặp được phủ
bởi 14 LP order (mỗi mệnh giá một order), tổng cộng $3 \times 2 \times
14 = 84$ LP order luôn được duy trì trong orderbook.

\subsubsection*{Pre-approval và refill loop}

Khi khởi động, Liquidity Service thực hiện hai pha:

\begin{enumerate}
  \item \textbf{Pre-approval.} Mỗi ví LP gọi \texttt{approve} lên
    contract USDC trên chuỗi nguồn của mình, cấp quyền cho hợp đồng
    HTLC chi tiêu một lượng lớn USDC (mặc định
    \texttt{type(uint256).max}). Pre-approval một lần duy nhất giúp
    các giao dịch \texttt{newOrder} sau đó không phải chèn thêm
    \texttt{approve}, tiết kiệm gas đáng kể. Nếu pre-approval thất
    bại, tiến trình crash để vận hành viên xử lý — không tự động
    fallback sang lịch trình chậm hơn.
  \item \textbf{Seed initial.} Liquidity Service bơm tất cả 84 LP
    order vào OrderStore. Mỗi LP order là một \texttt{IntentOrder}
    bình thường với \texttt{userAddress} là địa chỉ ví LP — nghĩa là
    với matching engine, nó không khác gì một intent người dùng.
\end{enumerate}

Sau đó, một \emph{refill loop} chạy với chu kỳ
\texttt{LP\_REFILL\_INTERVAL\_MS} (mặc định 10\,s). Mỗi vòng:

\begin{itemize}
  \item Quét các LP order đang được active map; với những vị trí mà
    LP order đã chuyển sang \texttt{MATCHED} hoặc \texttt{EXPIRED},
    tạo lại một LP order mới cùng mệnh giá để duy trì độ phủ.
  \item Ghi lại quan hệ "key (chuỗi-mệnh giá) $\rightarrow$ orderId hiện tại"
    trong cấu trúc \texttt{activeOrders} nội bộ.
\end{itemize}

Refill loop là cơ chế đảm bảo \emph{tính tự lành}: khi một LP order
được tiêu thụ bởi một cycle, một order thay thế xuất hiện trong
chậm nhất 10\,s, không cần sự can thiệp của vận hành viên.

\subsubsection*{Lọc LP-vs-LP để tránh self-match}

Một sự cố tiềm năng đáng chú ý: vì LP order trên chuỗi nguồn $X$ có
đích $Y$ và LP order trên chuỗi nguồn $Y$ có đích $X$, nếu được đưa
trực tiếp vào solver, hai chuỗi LP order \emph{ngược chiều} có thể
tạo thành cycle với nhau, dẫn đến việc các LP "tự match" lẫn nhau
mà không có người dùng nào hưởng lợi (chỉ tốn gas).

Liquidity Service xử lý điểm này tại \texttt{tick()} của chính nó,
không trong solver. Mỗi tick:

\begin{enumerate}
  \item Liệt kê các cặp chuỗi $(s, d)$ có \emph{user order} đang
    chờ.
  \item Khi đưa LP order vào snapshot cho matching engine, chỉ giữ
    lại các LP order có \emph{cặp chuỗi đối ứng} $(d, s)$ trùng với
    một cặp $(s, d)$ trong tập trên — nghĩa là các LP order có thể
    đóng vai trò "đối tác" cho user order đang sống.
  \item Các LP order trên cặp chuỗi không có user order tương ứng
    bị \emph{ẩn} khỏi solver trong tick đó. Chúng vẫn nằm trong
    OrderStore — chỉ là không tham gia matching cycle này.
\end{enumerate}

Cách lọc đẩy lên service-level (thay vì phải sửa solver) giữ cho
solver \emph{không biết khái niệm "LP"} — solver chỉ thấy một đồ
thị có hướng. Đây là một quyết định thiết kế quan trọng theo nguyên
tắc separation of concerns.

\subsubsection*{Chọn LP order cho một user intent}

Hàm \texttt{selectLpOrdersForUsers} (trong \texttt{lpSelector.ts})
nhận tập user intent và tập LP order, trả về tập LP order \emph{con}
tối thiểu cần thiết để phủ chính xác mỗi user amount. Vì LP order có
mệnh giá là luỹ thừa của 2 và mỗi mệnh giá xuất hiện đúng một lần
trên một cặp chuỗi, việc chọn quy về \emph{biểu diễn nhị phân của
$a$}: với mỗi bit 1 ở vị trí $i$, chọn LP order có mệnh giá $2^i$.

Ví dụ với $a = 700$ USDC:
$
700 = 512 + 128 + 32 + 16 + 8 + 4 = 2^9 + 2^7 + 2^5 + 2^4 + 2^3 + 2^2,
$
sáu LP order được chọn — phủ chính xác $700$ USDC, không thừa
không thiếu. Tính chất này khớp tốt với mô hình HTLC multi-fill: một
order consolidated trên chuỗi đích sẽ chứa sáu fill, mỗi fill ứng
với một LP order.

\subsubsection*{Liên kết với Matching Engine}

LP order về mặt cú pháp không khác gì user intent, nên solver xử lý
hỗn hợp LP + user trong cùng một đồ thị. Khác biệt chỉ ở:

\begin{itemize}
  \item LP order có \texttt{userAddress} là ví LP — Orchestrator nhận
    diện và dùng signer tương ứng (ví LP) khi tạo HTLC trên chuỗi
    nguồn.
  \item LP order \emph{không kèm} \texttt{incentiveFee}, do trọng số
    cạnh trong solver bằng đúng mệnh giá. Điều này làm các cạnh LP
    có trọng số đều và thấp hơn các intent người dùng có
    incentive — solver có khuynh hướng chọn user intent vào cycle
    trước, đúng với mục tiêu kinh tế.
\end{itemize}
